\section{Introduction}
\label{sec:intro}

\subsection{Longevity risk is priced off the interval}

A life annuity is a bet on a distribution. The best-estimate reserve of an
annuity book is a functional of the central mortality projection, but
everything that sits on top of it---the risk margin, the longevity component
of solvency capital, the price at which a longevity swap or a bulk annuity
transfer clears---is a functional of the projection's \emph{spread}. Under
the Solvency~II standard formula the longevity stress is a permanent
20\% fall in mortality rates \citep{eiopa2015delegated}; under an internal
model it is the 99.5\% one-year quantile of the change in own funds, which a
stochastic mortality model must supply from its own predictive distribution.
Either way the quantity carried on the balance sheet is a tail of a forecast
distribution, and its adequacy rests on one property: that a nominal 95\%
interval actually contains the realised outcome about 95\% of the time. An
interval that is too narrow under-reserves; one that is too wide over-prices
the hedge. Neither error is visible in a point-accuracy comparison.

Stochastic mortality models have produced intervals from the start.
\citet{lee1992modeling} forecast the period index $\kappa_t$ by a random walk
with drift and reported prediction intervals for life expectancy; the
Cairns--Blake--Dowd programme built a family of models expressly to generate
density forecasts \citep{cairns2006two,cairns2009quantitative,cairns2011mortality};
the actuarial toolkit \pkg{StMoMo} \citep{villegas2018stmomo} exposes bootstrap
and simulation intervals for all of them. The neural extensions of
Lee--Carter that followed
\citep{richman2021neural,nigri2019deep,perla2021time,scognamiglio2022calibrating}
improved point accuracy across many populations, and several now emit
intervals of their own
\citep{marino2023neural,miyata2022extending,schnurch2022point}. Whether
those intervals are \emph{calibrated} is a separate question from whether
they exist, and it is the question that decides what they are worth to an
actuary.

\subsection{The break that makes the question answerable}

Calibration is testable only where the truth is eventually observed, and the
test is informative only where the truth is allowed to surprise the model.
From the mid-1950s until 2019 the mortality surfaces of the high-income
populations in the Human Mortality Database (HMD) improved smoothly enough
that almost any trend extrapolation looked calibrated in sample. The years
2020--2024 broke that surface: period life expectancy fell in 2020 and 2021
in most HMD countries \citep{aburto2022quantifying}
\todo{verify magnitude and country coverage against Aburto et al.\ (2022)
before quoting any figure}, and the subsequent recovery was uneven by country
and by age. Any prediction interval constructed from data ending in 2019 was,
for the first time in the modern record, put to a genuine out-of-sample test
by reality.

The result of that test has not been scored. The closest study,
\citet{barigou2023bayesian}, has an exemplary evaluation
design---leave-future-out validation, log score and CRPS---but its pandemic
analysis is a simulation: French male deaths perturbed by ``two years of
excess mortality followed by one year of lower mortality'', because in 2021
real post-shock data did not exist. They exist now. The HMD vintage of
15~June~2026 (Methods Protocol v6) is the first with final single-age data
through 2024 for twenty populations. This paper uses it.

\subsection{What the literature measures}

We searched the full text of the mortality-forecasting library assembled for
this study (about twenty modelling papers; the term counts and quotations are
recorded in the project repository) for what each paper actually reports.
The counts below are over that library; where wider reading changed them,
the correction is stated.

\begin{itemize}
\item \emph{Point accuracy is the default and often the only metric.}
  \citet{richman2021neural}, the multi-population neural Lee--Carter that
  later work builds on, is fit to every HMD country and evaluated on mean
  squared error alone: the term ``MSE'' appears twenty-eight times in the
  full text; ``coverage'', ``CRPS'' and ``log score'' appear zero times.
  \citet{scognamiglio2022calibrating}, \citet{nigri2019deep},
  \citet{petnehazi2019mortality} and \citet{perla2021time} have the same
  profile.
\item \emph{Intervals are drawn but not checked.} \citet{miyata2022extending}
  advertise ``the ability to forecast with confidence intervals'', plot
  97.5\% intervals for the latent factor and for mortality at ages 30--80
  (their Section~6.4 and Figures~7--8), and then compare against Lee--Carter
  by MSE only (their Table~3). The paper never asks whether observed
  mortality falls inside the intervals it draws. The same authors note, of
  \citet{schnurch2022point}, that the intervals there are ``not based on the
  endogenous randomness \ldots\ but on the exogenous randomness driven by the
  random seed for the NN'': variance across seeds, which measures nothing
  about predictive coverage.
  \todo{Schn\"urch and Korn (2022) not yet read; confirm this characterisation
  and whether they tabulate coverage}
\item \emph{Where coverage is reported, it is reported without a sharpness
  penalty.} \citet{marino2023neural} is, to our knowledge, the one neural
  Lee--Carter paper that computes interval coverage. It reports the
  prediction-interval coverage probability (PICP) and the mean interval width
  (MPIW) as two separate indicators and treats $\mathrm{PICP}=1$ as the best
  outcome: the LSTM ``offers always a greater probability coverage, in most
  cases due to the PI width''. An interval wide enough to contain everything
  scores~$1$; without a width-penalised criterion, ``wider'' and ``better''
  are indistinguishable. Buried in the same paper, and not followed up, is
  the observation that the classical ARIMA intervals attain empirical
  coverage ``around 50\%'' and, for Spanish males, ``stable around
  33\%''---nominal 95\% Lee--Carter intervals covering a third to a half of
  realisations. The test window (2001--2018, three countries) contains no
  structural break. The most recent competitor, the gradient-boosted
  multi-population model of \citet{miao2025gradient}, repeats the pattern one
  level up: marginal coverages of 90\%, 94\% and 96\% against a nominal 95\%
  are reported and the closest is declared ``the most properly calibrated'',
  again with no proper score and no width penalty alongside.
\item \emph{Ex-post density evaluation exists, but predates both neural
  models and the break.} \citet{cairns2011mortality} assess the density
  forecasts of six stochastic models \emph{ex ante}, by the biological
  reasonableness of their fan charts; there is no PIT, no coverage
  computation, no proper score and no held-out evaluation.
  \citet{dowd2010backtesting} do backtest multi-period-ahead density
  forecasts ex post, for six classical models on English and Welsh males.
  \todo{VERIFY against Dowd et al.\ (2010b) once read: which backtests, whether
  hit rates or PIT values are tabulated, horizons, ages, and the verdict}
  That work is the direct ancestor of the present audit; it is confined to
  one population, to models with native intervals, and to a period without a
  break.
\item \emph{Absent from mortality forecasting, so far as we can find:} PIT or
  rank histograms; conformal prediction; joint (path) coverage over horizons;
  and, with the single exception of \citet{shang2018model}, formal
  forecast-comparison testing (Diebold--Mariano, model confidence sets).
  \todo{Crossref/Scholar search for ``conformal'' + ``mortality'' before
  submission; the library sweep was arXiv-biased}
\end{itemize}

The principle that the literature's coverage practice violates is old and
precise: a probabilistic forecaster should \emph{maximise the sharpness of
the predictive distributions subject to calibration}
\citep{gneiting2007probabilistic}. Coverage without width is not
calibration; width without coverage is not sharpness; and neither is a
proper score \citep{gneiting2007strictly}. The audit below scores every
forecast by both, and by the proper scores that combine them.

\subsection{This study}

The question is whether the prediction intervals of stochastic,
machine-learning and neural mortality forecasting models attain their
nominal coverage when the evaluation window crosses the COVID-19 break, and
whether the answer depends on the model family, on the uncertainty
mechanism, or on both. The last clause is the design. Existing comparisons
vary the model and inherit whatever uncertainty machinery each model happens
to ship with, so a coverage failure cannot be attributed: is the
architecture wrong about the future, or is the interval construction wrong
about the architecture? We treat \emph{model family} and \emph{uncertainty
mechanism} as separate crossed factors---ten families, seven mechanisms,
fifty admissible cells---so that the same architecture is audited under
several mechanisms and the same mechanism across several architectures.
That separation is what makes this a study rather than a bake-off.

Three regimes separate ordinary forecast error from break-induced failure: a
\emph{stable} control (expanding origins 1990--2014, every test year before
2020); the \emph{shift} regime (train to 2019, test 2020--2024); and a
\emph{placebo} break (train to 1913, test 1914--1922, spanning the First
World War and the 1918 influenza pandemic) that asks whether whatever we
find is specific to COVID-19 or is simply what intervals do at any break.

Five hypotheses were fixed and hash-stamped before any model touched real
data. They are reproduced verbatim from the pre-registration.

\begin{hypothesis}[H1]
Model rankings by RMSE and by CRPS disagree (rank correlation $<1$; at least
one top-3 inversion) in the shift regime.
\end{hypothesis}
\begin{hypothesis}[H2]
Nominal 95\% intervals under-cover in the shift regime for every model
family; the shortfall exceeds the stable-regime shortfall.
\end{hypothesis}
\begin{hypothesis}[H3]
Joint path coverage over $h = 1,\dots,5$ is materially below marginal
coverage at the same nominal level for every family (reported gap; test:
paired difference $>0$).
\end{hypothesis}
\begin{hypothesis}[H4]
Coverage failure is non-uniform in age: shift-regime coverage at ages 65--99
is lower than at ages 25--64 for every family.
\end{hypothesis}
\begin{hypothesis}[H5]
Interval miscalibration propagates to derived quantities: empirical coverage
of 95\% intervals for $e_{65}$ and annuity factor $\annuity$ (2\%) departs
from nominal by more than the corresponding stable-regime departure.
\end{hypothesis}

No directional claim is registered about neural-versus-classical ordering;
the crossed design estimates it.

\subsection{Contributions}

This paper proposes no new mortality model. Its contributions are:
\begin{enumerate}
\item \emph{The audit.} The first evaluation of probabilistic mortality
  forecasts---classical, neural and conformal---against realised 2020--2024
  mortality across twenty populations, with every interval constructed
  strictly from pre-2020 data.
\item \emph{The crossed design.} Model family and uncertainty mechanism as
  separate factors, so that a coverage failure can be attributed to the
  architecture, to the mechanism, or to their interaction. Claims are stated
  as contrasts within admissible sub-grids, never as main effects over a
  ragged grid.
\item \emph{Resolution of the failure.} By horizon (marginal versus joint
  path coverage, \Hthree), by age (\Hfour) and in economic terms (life
  expectancy and annuity-factor intervals, \Hfive), with the 1914--1922
  placebo as a second event.
\item \emph{The protocol.} A pre-registered, synthetic-truth-validated,
  oracle-parity-checked evaluation harness in which every model emits
  predictive samples through one interface and is scored through one code
  path, with cluster-bootstrap forecast-comparison inference that respects
  the fact that a pandemic is one shock, not twenty. The harness is reusable
  for any future break.
\end{enumerate}

Section~\ref{sec:related} places the audit in the literature.
Section~\ref{sec:data} describes the data and the conventions that make
Poisson scoring of fractional deaths well defined. Section~\ref{sec:design}
sets out the regimes, the grid, the interface and the validation gates.
Section~\ref{sec:metrics} defines every metric. Sections~\ref{sec:results}
and~\ref{sec:actuarial} are the results and their actuarial reading; both
are placeholders in this draft. Section~\ref{sec:discussion} pre-commits the
limitations and Section~\ref{sec:conclusion} concludes.
